% Patch record for C-19 -- table/figure layout cleanup (Table 1, Table 10, Figs. 43-46, 55-58, 59)
% Author-requested visual fixes, no data/claim changes. See PROGRESS.md C-19 row for full rationale.
% Raw diff below (introduction.tex + results.tex):

diff --git i/sections/introduction.tex w/sections/introduction.tex
index c29e8d4..fa06baf 100644
--- i/sections/introduction.tex
+++ w/sections/introduction.tex
@@ -13,7 +13,8 @@ A critical observation emerges from this body of work: in all published wheeled-
 \begin{table}[t]
 \centering
 \caption{Horizontal comparison of bio-inspired and engineering approaches to locomotion-mode action selection.}
-\begin{tabular}{|p{2.6cm}|p{2.9cm}|p{2.3cm}|p{1.6cm}|p{2.6cm}|p{2.4cm}|}
+\footnotesize
+\begin{tabular}{|p{2.1cm}|p{2.6cm}|p{1.9cm}|p{1.4cm}|p{2.4cm}|p{2.0cm}|}
 \hline
 \textbf{Approach} & \textbf{Selection Mechanism} & \textbf{Role in Control Stack} & \textbf{Biological Basis} & \textbf{Interpretability / Training} & \textbf{Representative Work} \\ \hline
 Hand-coded threshold logic & Fixed if-then rules on sensor thresholds & Arbitration & None & Fully transparent; manually tuned, no training & \citet{Kashiri2019CENTAURO} \\ \hline
diff --git i/sections/results.tex w/sections/results.tex
index 89fbc75..5e2e654 100644
--- i/sections/results.tex
+++ w/sections/results.tex
@@ -544,41 +544,40 @@ During this movement, a leftward correction occurs to center the stimulus, refle
 
 To complement the qualitative description of the terrain-classification module (Section~\ref{ssec:Gait_Decision_Module}), the MLP (150–80 hidden units, \revblue{LeakyReLU with negative slope $0.01$}, Adam optimizer) was evaluated on the IMU dataset ($n = 600$ samples, 135-dimensional input, binary labels: 0 = flat terrain, 1 = inclined terrain) using a stratified 80/20 train–test split and 5-fold stratified cross-validation.
 
-As illustrated in Figure~\ref{fig:confusion_matrix}, the final model produced a confusion matrix on the held-out test set. Across the 5-fold cross-validation, the model achieved an accuracy of $0.767 \pm 0.028$, precision of $0.764 \pm 0.024$, recall of $0.779 \pm 0.048$, and an F1-score of $0.771 \pm 0.033$. On the held-out test set ($n = 120$), the model reached an accuracy of 0.800, precision of 0.747, recall of 0.918, F1-score of 0.824, and an ROC-AUC of 0.866, confirming that the network reliably discriminates between flat and inclined terrain, with a recall bias toward the inclined-terrain class that minimises the safety-critical risk of missed inclines.
+\revblue{As illustrated in Figure~\ref{fig:mlp_eval_grid}a}, the final model produced a confusion matrix on the held-out test set. Across the 5-fold cross-validation, the model achieved an accuracy of $0.767 \pm 0.028$, precision of $0.764 \pm 0.024$, recall of $0.779 \pm 0.048$, and an F1-score of $0.771 \pm 0.033$. On the held-out test set ($n = 120$), the model reached an accuracy of 0.800, precision of 0.747, recall of 0.918, F1-score of 0.824, and an ROC-AUC of 0.866, confirming that the network reliably discriminates between flat and inclined terrain, with a recall bias toward the inclined-terrain class that minimises the safety-critical risk of missed inclines. \revblue{The classifier's discrimination capacity is further supported by the ROC curve (Figure~\ref{fig:mlp_eval_grid}b), where the curve remains well above the chance diagonal across the operating range, yielding an AUC of 0.866.}
 
-\begin{figure}[htbp]
+\revblue{Figure~\ref{fig:mlp_eval_grid}c illustrates the training convergence behaviour}, showing the evolution of the learning loss across epochs and confirming stable optimisation of the selected configuration. To justify the choice of architecture, an ablation study compared the proposed 150–80 configuration against simpler (50, 100) and deeper (200, 100, 50) topologies. \revblue{As shown in Figure~\ref{fig:mlp_eval_grid}d}, the 150–80 configuration achieved the highest F1-score among all tested architectures ($\text{F1} = 0.771$), outperforming the shallower (50: $\text{F1} = 0.747$; 100: $\text{F1} = 0.766$) and deeper (200-100-50: $\text{F1} = 0.766$) variants. This confirms that 150–80 represents the most favourable trade-off between model capacity and generalisation for this classification task, while remaining computationally inexpensive for onboard deployment.
+
+\begin{figure*}[t]
     \centering
-    \includegraphics[width=0.8\textwidth]{fig1_confusion_matrix.png}
-    \caption{Confusion matrix of the final model (150-80) evaluated on the test set.}
-    \label{fig:confusion_matrix}
-\end{figure}
+    \begin{subfigure}[b]{0.48\textwidth}
+        \centering
+        \includegraphics[width=\linewidth]{fig1_confusion_matrix.png}
+        \caption{Confusion matrix (test set).}
+    \end{subfigure}
+    \hfill
+    \begin{subfigure}[b]{0.48\textwidth}
+        \centering
+        \includegraphics[width=\linewidth]{fig2_roc_curve.png}
+        \caption{ROC curve (AUC = 0.866).}
+    \end{subfigure}
 
-The classifier's discrimination capacity is further supported by the ROC curve shown in Figure~\ref{fig:roc_curve}, where the curve remains well above the chance diagonal across the operating range, yielding an AUC of 0.866.
+    \vspace{0.6em}
 
-\begin{figure}[htbp]
-    \centering
-    \includegraphics[width=0.8\textwidth]{fig2_roc_curve.png}
-    \caption{Receiver Operating Characteristic (ROC) curve of the classifier (AUC = 0.866).}
-    \label{fig:roc_curve}
-\end{figure}
-
-Figure~\ref{fig:training_loss} illustrates the training convergence behaviour, showing the evolution of the learning loss across epochs and confirming stable optimisation of the selected configuration.
-
-\begin{figure}[htbp]
-    \centering
-    \includegraphics[width=0.8\textwidth]{fig3_training_loss.png}
-    \caption{Training convergence curve displaying learning loss versus epoch.}
-    \label{fig:training_loss}
-\end{figure}
-
-To justify the choice of architecture, an ablation study compared the proposed 150–80 configuration against simpler (50, 100) and deeper (200, 100, 50) topologies. As shown in Figure~\ref{fig:architecture_ablation}, the 150–80 configuration achieved the highest F1-score among all tested architectures ($\text{F1} = 0.771$), outperforming the shallower (50: $\text{F1} = 0.747$; 100: $\text{F1} = 0.766$) and deeper (200-100-50: $\text{F1} = 0.766$) variants. This confirms that 150–80 represents the most favourable trade-off between model capacity and generalisation for this classification task, while remaining computationally inexpensive for onboard deployment.
-
-\begin{figure}[t]
-    \centering
-    \includegraphics[width=\textwidth]{fig4_architecture_ablation.png}
-    \caption{Architecture comparison (accuracy, precision, recall, F1).}
-    \label{fig:architecture_ablation}
-\end{figure}
+    \begin{subfigure}[b]{0.48\textwidth}
+        \centering
+        \includegraphics[width=\linewidth]{fig3_training_loss.png}
+        \caption{Training convergence.}
+    \end{subfigure}
+    \hfill
+    \begin{subfigure}[b]{0.48\textwidth}
+        \centering
+        \includegraphics[width=\linewidth]{fig4_architecture_ablation.png}
+        \caption{Architecture ablation (F1-score).}
+    \end{subfigure}
+    \caption{\revblue{MLP-based terrain classifier evaluation, final configuration (150--80 hidden units). (a) Confusion matrix on the held-out test set ($n=120$). (b) ROC curve (AUC = 0.866). (c) Training-loss convergence across epochs. (d) F1-score comparison across the four tested hidden-layer architectures, motivating the 150--80 choice.}}
+    \label{fig:mlp_eval_grid}
+\end{figure*}
 
 The MLP was further benchmarked against alternative classifiers (Logistic Regression, SVM with RBF kernel, and Random Forest) under the same cross-validation protocol. Logistic Regression underperformed ($\text{F1} = 0.692$), reflecting the non-linear separability of the IMU features, while SVM-RBF (\revblue{$\text{F1} = 0.857$}) and Random Forest (\revblue{$\text{F1} = 0.947$}) achieved higher raw scores (Figure~\ref{fig:classifier_tradeoff}a); the rationale for nonetheless selecting the MLP is given below, alongside the per-sample latency comparison.
 
@@ -673,7 +672,7 @@ To further consolidate the physical validation under complex, multi-section topo
 \centering
 \caption{Gait-mode transitions across the six-frame compound-terrain sequence (Figure~\ref{fig:real_variable_topography_sequence}), replicated across 15 continuous trials.}
 \label{tab:terrain_sequence}
-\begin{tabularx}{\textwidth}{c l l X}
+\begin{tabularx}{\textwidth}{>{\centering\arraybackslash}p{0.08\textwidth} >{\raggedright\arraybackslash}p{0.19\textwidth} >{\raggedright\arraybackslash}p{0.19\textwidth} >{\raggedright\arraybackslash}X}
 \toprule
 \textbf{Frame} & \textbf{Terrain segment} & \textbf{Mode} & \textbf{Trigger / mechanism} \\
 \midrule
@@ -746,14 +745,14 @@ This subsection presents a rigorous quantitative analysis of the proposed bio-in
 
 \begin{figure*}[t]
   \centering
-  \begin{subfigure}[b]{0.32\textwidth}
+  \begin{subfigure}[b]{0.48\textwidth}
     \centering
     \includegraphics[width=\linewidth]{Appetitive_Real_Plot_1_Switching_Delay_ECDF.png}
     \caption{Switching delay ECDF.}
     \label{fig:appetitive_ecdf_real}
   \end{subfigure}
   \hfill
-  \begin{subfigure}[b]{0.32\textwidth}
+  \begin{subfigure}[b]{0.48\textwidth}
     \centering
     \includegraphics[width=\linewidth]{Appetitive_Real_Plot_3_Temporal_Dynamics_PitchRollRMS.png}
     \caption{Pitch/Roll RMS attitude dynamics.}
@@ -769,14 +768,14 @@ This subsection presents a rigorous quantitative analysis of the proposed bio-in
 
 \begin{figure*}[t]
   \centering
-  \begin{subfigure}[b]{0.32\textwidth}
+  \begin{subfigure}[b]{0.48\textwidth}
     \centering
     \includegraphics[width=\linewidth]{Aversive_Real_Plot_1_Switching_Delay_ECDF.png}
     \caption{Switching delay ECDF.}
     \label{fig:aversive_ecdf_real}
   \end{subfigure}
   \hfill
-  \begin{subfigure}[b]{0.32\textwidth}
+  \begin{subfigure}[b]{0.48\textwidth}
     \centering
     \includegraphics[width=\linewidth]{Aversive_Real_Plot_3_Temporal_Dynamics_PitchRollRMS.png}
     \caption{Pitch/Roll RMS attitude dynamics.}
@@ -792,14 +791,14 @@ This subsection presents a rigorous quantitative analysis of the proposed bio-in
 
 \begin{figure*}[t]
   \centering
-  \begin{subfigure}[b]{0.32\textwidth}
+  \begin{subfigure}[b]{0.48\textwidth}
     \centering
     \includegraphics[width=\linewidth]{Complex_Real_Plot_1_Switching_Delay_ECDF.png}
     \caption{Switching delay ECDF.}
     \label{fig:complex_ecdf_real}
   \end{subfigure}
   \hfill
-  \begin{subfigure}[b]{0.32\textwidth}
+  \begin{subfigure}[b]{0.48\textwidth}
     \centering
     \includegraphics[width=\linewidth]{Complex_Real_Plot_3_Temporal_Dynamics_PitchRollRMS.png}
     \caption{Pitch/Roll RMS attitude dynamics.}
@@ -815,14 +814,14 @@ This subsection presents a rigorous quantitative analysis of the proposed bio-in
 
 \begin{figure*}[t]
   \centering
-  \begin{subfigure}[b]{0.32\textwidth}
+  \begin{subfigure}[b]{0.48\textwidth}
     \centering
     \includegraphics[width=\linewidth]{Obstacle_Real_Plot_1_Switching_Delay_ECDF.png}
     \caption{Switching delay ECDF.}
     \label{fig:obstacle_ecdf_real}
   \end{subfigure}
   \hfill
-  \begin{subfigure}[b]{0.32\textwidth}
+  \begin{subfigure}[b]{0.48\textwidth}
     \centering
     \includegraphics[width=\linewidth]{Obstacle_Real_Plot_3_Temporal_Dynamics_PitchRollRMS.png}
     \caption{Pitch/Roll RMS attitude dynamics.}
@@ -854,7 +853,7 @@ Obstacle   & ${\approx}\,0.0\,\text{s}$ (instantaneous) & ${\approx}\,1100$, tra
 
 \begin{figure}[t]
   \centering
-  \includegraphics[width=0.75\textwidth]{FigReal_Morphological_Transition_Energy.png}
+  \includegraphics[width=0.5\textwidth]{FigReal_Morphological_Transition_Energy.png}
   \caption{\revblue{Morphological transition energy per direction, derived from $|I_{\text{current}}| \times 7.2\,\text{V}$ (fixed NiMH pack voltage) integrated over the pre-calibrated electromechanical coupling window (2.92\,s / 1.56\,s, matching Fig.~\ref{fig:FourPanelLatencies}), across $N=16$ (Quadrupedal $\rightarrow$ Differential Rolling) and $N=19$ (Differential Rolling $\rightarrow$ Quadrupedal) individually detected structural-command transitions. Error bars: $\pm 1$ SD; the between-direction difference is not statistically significant (Welch's $t \approx 1.18$).}}
   \label{fig:MorphologicalEnergy}
 \end{figure}
